\newcommand{\figuraGeometriaTierraLuna}{%
\begin{figure}[H]
    \centering
    \begin{tikzpicture}[>=Latex,font=\small,line cap=round,line join=round]
        \coordinate (O) at (1.8,2.2);
        \coordinate (C) at (2.62,2.2);
        \coordinate (L) at (10.6,2.2);

        \draw[densely dashed,black!45] (0.15,2.2)--(11.35,2.2);
        \draw[->,thick] (O)--(11.55,2.2)
            node[above] {$+z$};

        \fill[blue!18] (O) circle (1.18);
        \draw[very thick,blue!65!black] (O) circle (1.18);
        \draw[blue!45!white,thick] (1.15,2.55)
            .. controls (1.45,2.92) and (1.78,2.82) .. (2.05,3.22);
        \draw[blue!45!white,thick] (1.18,1.58)
            .. controls (1.55,1.35) and (1.72,1.75) .. (2.30,1.50);
        \node[font=\bfseries,blue!65!black] at (1.8,0.67) {Tierra};

        \fill[gray!38] (L) circle (0.44);
        \draw[very thick,black!65] (L) circle (0.44);
        \fill[gray!55] (10.45,2.34) circle (0.08);
        \fill[gray!55] (10.78,2.09) circle (0.06);
        \node[font=\bfseries] at (10.6,1.42) {Luna};

        \fill[black] (O) circle (0.055);
        \node[anchor=north east] at ($(O)+(-0.04,-0.05)$) {$O$};
        \fill[red!75!black] (C) circle (0.075);
        \node[red!75!black,anchor=south] at ($(C)+(0,0.10)$)
            {$C$};

        \draw[<->,thick] (1.8,0.18)--(10.6,0.18)
            node[midway,fill=white,inner sep=1pt] {$D$};
        \draw[densely dashed] (1.8,0.35)--(1.8,1.02);
        \draw[densely dashed] (10.6,0.35)--(10.6,1.72);

        \draw[->,thick,blue!65!black] (O)--($(O)+(0,1.18)$)
            node[midway,left] {$a$};
        \draw[red!75!black,thick] (C)--(3.38,3.32);
        \node[red!75!black,anchor=west] at (3.42,3.32)
            {centro de masa común};

        \node[align=center,text=black!70] at (6.3,1.15)
            {Esquema sin escala\\$\vec{x}_0=D\hat z$};
    \end{tikzpicture}
    \caption{Geometría simplificada del sistema Tierra--Luna. El centro de
    masa común $C$ se encuentra dentro de la Tierra; el origen $O$ se toma en
    el centro terrestre y el eje $z$ apunta hacia la Luna.}
    \label{fig:geometria-tierra-luna}
\end{figure}%
}